\section{Two Convergence Rules}
\label{sec:convergence}

The four failure modes in Section~\ref{sec:implementation} concern correctness.
Two additional design rules concern convergence quality: they do not prevent learning
from starting but determine whether it converges reliably across different task
geometries and random seeds.
Both rules emerged from Study~1 and are presented here as transferable design
principles before the experimental studies.

\subsection{Lengthscale-to-Range Scaling Rule}
\label{sec:convergence:lengthscale}

\paragraph{The problem.}
The RBF policy (Eq.~\ref{eq:policy}) uses a squared-exponential kernel with
lengthscale $\ell_s$.
Define the \emph{RBF sensitivity} between two target positions at Euclidean
distance $d$ as
\begin{equation}
  S(d,\, \ell_s) = \exp\!\left( -\frac{d^2}{2\,\ell_s^2} \right).
  \label{eq:rbf_sensitivity}
\end{equation}
$S \to 1$ means the two RBF values are indistinguishable, and the policy outputs the
same speed for both targets.
$S \to 0$ means the targets are well-separated and the policy can assign independent
speeds.

The paper's default $\ell_s = 1.0$\,m is calibrated to a target range of
$\approx 1.0$\,m (Config~A), where $S(1.0,\, 1.0) = e^{-0.5} \approx 0.607$ —
distinguishable.
For Config~E (target range $\Delta = 0.35$\,m), the same lengthscale gives
$S(0.35,\, 1.0) = e^{-0.061} \approx 0.941$ — effectively identical.
All 250 RBF centres output the same value for every target in the workspace, the
gradient is near-zero everywhere, and the policy degenerates to a constant-speed output.
We call this \emph{RBF blindness}.

\paragraph{The rule.}
Setting $\ell_s \approx 0.15 \times \Delta_{\mathrm{range}}$ where
$\Delta_{\mathrm{range}} = \ell_M - \ell_{\min}$ is the width of the target workspace
yields $S \approx 0.066$, matching the discriminability level of the working baseline.
Formally:
\begin{equation}
  \ell_s^* \approx 0.15\;\Delta_{\mathrm{range}}.
  \label{eq:ls_rule}
\end{equation}
The derivation in Appendix~\ref{app:lengthscale} shows that $\ell_s = \Delta_{\mathrm{range}}$
exactly preserves Config~A's sensitivity; the empirical factor of 0.15 is conservative
and provides faster convergence when GP data are sparse.

\paragraph{Validation.}
Table~\ref{tab:sensitivity} verifies the rule across all height configurations.
Config~E with $\ell_s^* = 0.15$\,m achieves 10/10 hits immediately after this
correction; four prior iterations all failed with $\ell_s = 1.0$\,m despite correct
GP dynamics (Section~\ref{sec:study1:config_e}).

\begin{table}[t]
  \centering
  \caption{RBF sensitivity $S$ at target range $\Delta_{\mathrm{range}}$ for each
           configuration.
           Config~E with default $\ell_s = 1.0$\,m is near-unity; E-Strat5 with
           $\ell_s = 0.15$\,m recovers the working discriminability of Config~A.}
  \label{tab:sensitivity}
  \small
  \begin{tabular}{lcccc}
    \toprule
    Config & $\Delta_{\mathrm{range}}$ (m) & $\ell_s$ (m) & $S$ & Result \\
    \midrule
    A              & $\approx 1.00$ & 1.00 & 0.607 & 5/5 \\
    B-Strat        & $\approx 1.00$ & 1.00 & 0.607 & 10/10 \\
    C-Strat        & $\approx 1.00$ & 1.00 & 0.607 & 10/10 \\
    D-Strat        & $\approx 1.00$ & 1.00 & 0.607 & 10/10 \\
    E (default)    & 0.35           & 1.00 & 0.941 & 0/10 \textbf{FAIL} \\
    E-Strat5       & 0.35           & 0.15 & 0.066 & 10/10 \\
    \bottomrule
  \end{tabular}
\end{table}

\paragraph{Practical implication.}
Equation~\eqref{eq:ls_rule} is an \emph{a priori} design check: given the task geometry,
verify $S < 0.1$ at the extremes of the target workspace before running any experiment.
If $S > 0.5$, increase $N_b$ or reduce $\ell_s$.

\subsection{Stratified Exploration Rule}
\label{sec:convergence:stratified}

\paragraph{The problem.}
Random exploration with $N_{\mathrm{exp}} = 5$ throws draws velocities uniformly from
$[0, u_M]$.
An unlucky seed can draw four low-speed throws that all land well short of the target
zone, leaving the GP with no observations in the speed region needed to reach any target.
When the policy gradient subsequently pushes commanded speed toward $u_M$, the GP must
extrapolate beyond its training support — and it cannot.
Config~B with seed=1 experienced exactly this: five exploration throws produced landing
positions of 0.245\,m and 0.408\,m from the origin, while the closest target was at
$\approx 1.43$\,m.
The policy saturated at $u_M$ and all 10 evaluation throws missed.
The identical configuration with seed=2 achieved 10/10 hits, confirming that the
failure was a property of the seed, not the algorithm.

\paragraph{The fix: stratified band sampling.}
We introduce \texttt{Stratified\_Throwing\_Exploration}, which partitions the speed
interval $[0, u_M]$ into $N_{\mathrm{exp}}$ equal bands and samples throw $i$ uniformly
from band $i$:
\begin{equation}
  v_i \sim \mathcal{U}\!\left(
    \frac{(i-1)\,u_M}{N_{\mathrm{exp}}},\;
    \frac{i\,u_M}{N_{\mathrm{exp}}}
  \right), \quad i = 1, \ldots, N_{\mathrm{exp}}.
  \label{eq:stratified}
\end{equation}
This guarantees that the GP receives at least one observation in each fifth of the speed
range regardless of seed, with zero additional computational overhead.

\paragraph{Evidence.}
Applying stratified exploration to Config~B with seed=1 recovers 10/10 hits
(B-Strat, Table~\ref{tab:sensitivity}).
The cost trajectory at seed=1 with random exploration plateaus at 0.70--0.73 (max-speed
local minimum); with stratified exploration it falls to 0.024 within the first trial.
